\documentclass[parskip=half, american]{scrartcl}
\usepackage{booktabs}
\usepackage[margin=2.5cm]{geometry}
\usepackage[dvipsnames]{xcolor}
\usepackage[colorlinks=true,
			citecolor=blue,
			urlcolor=blue,
			linkcolor=blue,
			pdfborder={0 0 0}]
	{hyperref}
	
\newcommand{\AR}[1]
	{\color{PineGreen}AR: #1\color{black} \par }
	
\newcommand{\journal}{\emph{Fire}}


\title{\Large Remote sensing in rangeland fire ecology: Comparing imagery to measured fire behavior, and burn severity across prescribed burns and wildfires  }
\subtitle{Response to review}
\author{ }
\date{}

\pagenumbering{gobble}
\raggedright 

\begin{document}

\maketitle

\vspace{-2cm} 

\AR{Author Responses (AR) appear below in this green text.}

\section*{Reviewer 3}

This manuscript presents a study that addresses a critical knowledge gap in Rangeland Fire Ecology. I recommend publication after minor revisions, primarily to enhance the discussion of certain methodological nuances and to strengthen the comparison with existing literature.

\AR{I appreciate the reviewer's support and helpful comments.  }

In the abstract add a concise management takeaway, such as confirming $\Delta$NBR as a proxy for direct fire behavior measurements.

\AR{ Abstract includes ``By describing meaningful gradients in surface fire behavior in rangelands with $\Delta$NBR, even those without the capacity to measure fire behavior in the field can monitor prescribed fire effectiveness and incorporate burn severity in adaptive management plans.'' }

Briefly mention the ``two zones'' in the introduction to improve the transition to methods and results.

\AR{This has been added to the final paragraph of the Introduction.  }

Related Work: Add a brief sub-section (e.g., ``Synthesis and Knowledge Gap'') summarizing the state of the art and clearly showing how this study goes beyond previous work, especially for rangelands. 

\AR{I appreciate this point although it isn't clear where the reviewer wishes to see the sub-section added.
In the spirit of the comment I've attempted to revise relevant parts of the Introduction and Discussion accordingly.  }

4. Methodology

Methodology: Briefly justify the 15 cm thermocouple height and clarify how the ``flame temperature'' metric was derived from the time-series data.

\AR{Notes on these points have been added.  } 

Results: Add a table summarizing the mean and range (or standard deviation) of $\Delta$NBR, flame temperature, and rate of spread for prescribed burns and wildfires in each zone to show group variability.

\AR{Table added. }
	
Results Comparison with Related Works: Explicitly contrast the low-to-moderate severities observed here with the higher severities reported in forest fire studies (e.g., MTBS) to highlight the unique fire effects in rangelands.

\AR{This is a useful comment, I appreciate it. Established thresholds for fire severity classifications have been added to the relevant figure to contextualize the data, and this point has been added to the Discussion. }


Conclusion: Briefly reiterate the partially supported hypothesis on soil temperature and add a one-sentence explanation, e.g., that soil heating in grass fires depends on residence time, not captured by post-fire severity indices.

\AR{This is not actually consistent with the paper. 
	Low values for post-fire severity indices reflect the fact that grass fires do not have long residence times.
	There is no evidence that were longer residence times to occur, that they would not be captured by post-fire severity indices.  }

References: Ensure all in-text citations match the reference list and verify that DOIs or URLs are correctly formatted.

\AR{I appreciate the comment.} 

English Writing: Ensure consistent italics for foreign phrases and correct minor typographical issues, such as spacing after citations.

\AR{I appreciate the comment.} 


\end{document}
